% Research fragment: all labels begin ten:. Standard amsmath/amsthm required.
% Cite AlternatingIncoming, QuarticIncoming, and GeneralHankel as the inspected project sources.
\section{Complete reconstruction for smaller tail order ten}
\label{ten:section}

The two latest incoming reconstruction reports leave a particularly
small but genuine gap.  The alternating-tail report proves the order-ten
prefix theorem outside $b\equiv962\pmod{2048}$, while the quartic
reconstruction report proves it for every $b\ge1920$
\cite{AlternatingIncoming,QuarticIncoming}.  We complete the family.  The new step is
a coefficientwise congruence with one additional bit of precision:
its only exceptional class is $b\equiv3010\pmod{4096}$, which lies
entirely inside the range of the second report's real-variable argument.
We reproduce that argument at order ten, with explicit constants, to
make the resulting nonvanishing proof complete.  The small determinant
normalization and its Pad\'e origin are credited to the source reports;
the modulus-$2^{17}$ congruence and the completed family are the present
contribution.  No claim of arithmetic independence of numerical special
values is involved.

\subsection{Statement and exact meaning of reconstruction}

Let
\[
 A_s(n)=\zeta(s,n+1)=\zeta(s)-H_n^{(s)},\qquad x_n=n+\tfrac12,
 \qquad
 D_{10,2b}(w)=\sum_{n\ge0}A_{10}(n)A_{2b}(n)x_n^{-w}.
\]
The series initially converges for $\operatorname{Re}w>-2b-7$.
Its centered Euler--Maclaurin continuation is regular at nonpositive
even integers.  Write $D_r=D_{10,2b}(-2r)$ and use the source's
triangular coordinates
\begin{align}
 C_j&=\sum_{r=0}^{j}\binom{2j+1}{2r}2^{-2(j-r)}D_r,
 \label{ten:C}\\
 D_r&=\sum_{j=0}^{r}\frac{\binom{2r}{2j}}{2j+1}
             B_{2r-2j}(1/2)C_j.
 \label{ten:D}
\end{align}
For $j\ge b$, each $C_j$ is a rational combination of
$Z_k=\zeta(2k)$, $0\le k<b$, with $Z_0=-1/2$.  The words
\emph{formal rank} refer to rational coefficient vectors in the
formal basis
\[
 C_0,\ldots,C_{b-1},Z_0,\ldots,Z_{b-1};
\]
they do not refer to the linear independence of the values after
evaluation.

Put $\mu_j=B_{2j}/(2j)!$ for $j\ge0$, and $\mu_j=0$ for $j<0$.
The late coordinate matrix is
\begin{equation}\label{ten:M}
 (M_{5,b})_{tk}=\mu_{t+k-4}
 \left\{\frac{(2b+2t-9)!}{(2b-2k)!}
 +\mathbf1_{k<5}\frac{(2t+1)!}{(10-2k)!}\right\},
 \qquad 0\le t,k<b.
\end{equation}
In particular,
$(C_b,\ldots,C_{2b-1})^{\mathsf T}
=M_{5,b}(Z_0,\ldots,Z_{b-1})^{\mathsf T}$.
Formula \eqref{ten:M} also defines the late coordinate row for
every integer $t\ge0$.  Its binomial form, useful without factorial
conventions, is
\[
 B_{2t+2k-8}\left\{
 \frac{\binom{2t+2}{10-2k}}{2t+2}
 +\frac{\binom{2b+2t-8}{2b-2k}}{2b+2t-8}\right\},
\]
where negative Bernoulli indices give zero and a summand is omitted
unless its lower binomial index lies between $2$ and its upper index.
These are the inherited finite Bernoulli moment formulas
\cite{GeneralHankel}.  The full first-$2b$ coordinate determinant is
$\det M_{5,b}/(4b-1)!!$.

\begin{theorem}[Complete order-ten optimal prefix]\label{ten:main}
For every integer $b\ge6$, $M_{5,b}$ is nonsingular.  Consequently,
for every even $q>10$, the $q$ moments
\[
 D_{10,q}(0),D_{10,q}(-2),\ldots,D_{10,q}(-2(q-1))
\]
have full formal rank $q$.  Every subsequent centered even moment
is an explicit rational linear combination of this prefix.
\end{theorem}

The conclusion now holds for every unequal even pair whose smaller
order is at most ten, using the previously established smaller-order
families.  It does not establish the unrestricted conjecture at
smaller orders twelve and above, nor the full gap-six family.

\subsection{A finite five-dimensional definition of the numerator}

The following normalization is retained from \cite{AlternatingIncoming}.
For $0\le i,\ell\le4$, define
\[
 n_i=b-5+i,\quad N_{i\ell}=2(n_i-\ell)+1,
 \quad R_\ell=2b-8+2\ell,
\]
\[
 F_{i\ell}=\frac1{(2\ell+1)!}
       \prod_{r=0}^{2\ell-1}(2n_i-r)(2n_i+3+r),
\]
and define Newton coefficients $d_{ih}$ by the finite identity
\begin{equation}\label{ten:Newton}
 (2n_i+2-y)(4n_i+2-y)^{\underline{2i}}
       =\sum_{h=0}^{2i+1}d_{ih}y^{\underline h}.
\end{equation}
Equivalently,
\[
 d_{ih}=\frac1{h!}\sum_{r=0}^{h}(-1)^{h-r}\binom hr
       (2n_i+2-r)(4n_i+2-r)^{\underline{2i}}.
\]
For independent variables $U_1,\ldots,U_9$ put
\begin{equation}\label{ten:calA}
 \mathcal A_{i\ell}=2F_{i\ell}+
 \frac{R_\ell!}{(2n_i+1)(2n_i+2)(2\ell+2)!N_{i\ell}!}
 \sum_{k=0}^{i}d_{i,2k+1}N_{i\ell}^{\underline{2k+1}}
                 U_{5-i+\ell+k}.
\end{equation}
All factorial ratios with $b$-dependent arguments have fixed index
differences and therefore mean finite rational products in $b$.

Set
\begin{align*}
 \Delta_{10}(b)={}&4096b^{12}-110592b^{11}+1236992b^{10}
 -7303680b^9+23705088b^8\\
 &-38140416b^7+14925056b^6+24490080b^5-4469984b^4\\
 &-15838992b^3-2571948b^2+4244400b+1488375,
\end{align*}
and define the finite rational expression
\begin{equation}\label{ten:Pdef}
 P_{10}(b;U)=138240(2b-7)(2b-5)\Delta_{10}(b)
                       \frac{\det\mathcal A}{\det F}.
\end{equation}
Exact coefficient arithmetic gives
\begin{align}\label{ten:Fdet}
 \det F={}&\frac{16384}{4465125}(b-3)(b-2)^2(b-1)
     (2b-7)(2b-5)^2(2b-3)\notag\\
 &\cdot(4b-15)(4b-13)(4b-11)^2(4b-9)^2(4b-7)^2
          (4b-5)(4b-3)\Delta_{10}(b).
\end{align}
The expression $P_{10}$ belongs to $\mathbb Z[b,U_1,\ldots,U_9]$.
It has $439$ monomials when grouped by the $U$ exponents, or
$10403$ monomials after expanding also in $b$.  Its coefficients
have degree at most $24$ in $b$.  Equations
\eqref{ten:Newton}--\eqref{ten:Pdef} specify it completely, without
interpolation or a dependence on an unspecified fitted expression.
The accompanying exact verifier independently reconstructs
\eqref{ten:Fdet}, clears the five row denominators of
\eqref{ten:calA}, sums the $120$ signed permutation products, and
compares every coefficient of \eqref{ten:Pdef} with the inherited
integer certificate.  Thus the reused certificate is independently
replayed as an algebraic identity.

For $b\ge10$, with $H=(\mu_{t+k-4})_{0\le t,k<b}$, define
\[
 \rho_{5,b}=\frac{\det M_{5,b}}
 {\displaystyle
  \frac{\prod_{t=0}^{b-1}(2b+2t-9)!}
       {\prod_{k=0}^{b-1}(2b-2k)!}\det H}.
\]
The source reduction is
\begin{equation}\label{ten:rho}
 \rho_{5,b}
 =\left.\frac{\det\mathcal A}{\det F}\right|_{U_j=2B_{2b-2j}}
 =\frac{P_{10}(b;2B_{2b-2},\ldots,2B_{2b-18})}
 {138240(2b-7)(2b-5)\Delta_{10}(b)}.
\end{equation}
Here $H$ is nonsingular.  One direct proof conjugates $H$ by
$\operatorname{diag}(1,-1,1,\ldots)$ and partitions after its first
three rows and columns.  The resulting matrix has blocks
\[
 \begin{pmatrix}ee^{\mathsf T}&B\\B^{\mathsf T}&-P\end{pmatrix},
 \qquad e=(0,0,1)^{\mathsf T},
\]
where $P$ is positive definite by the representation
$2\sum_{m\ge1}(2\pi m)^{-2s}p((2\pi m)^{-2})^2>0$ for
$s\ge1$ and a nonzero polynomial $p$.  The first two rows of $B$
have a nonsingular antitriangular minor.  Consequently
$ee^{\mathsf T}+BP^{-1}B^{\mathsf T}$ is positive definite.

For completeness, the exact nonzero multiplier behind
\eqref{ten:rho} can be obtained without assuming Pad\'e normality.
The monic polynomials
\[
 \pi_n(x)=\sum_{t=0}^{n}
 \frac{2(2n+2t+1)!(2n+1)!}
 {(4n+2)!(2n-2t+1)!(2t)!}x^t
\]
satisfy, for $\mathcal N(x^r)=\mu_{r+1}$,
\[
 \mathcal N(x^j\pi_n)=0\ (j<n),\qquad
 \mathcal N(x^n\pi_n)=h_n,
 \quad h_n=\frac{((2n+1)!)^2}{(4n+2)!(4n+3)!}.
\]
These finite identities follow by coefficient comparison in the
diagonal Pad\'e formula for $e^t$, as in \cite{QuarticIncoming}; they
can also be checked directly using the Bernoulli generating function.
The same coefficients give
\[
 \sum_{t=0}^{n}[x^t]\pi_n(x)\,\mu_{t-\ell}
 =c_n f_\ell(n),\qquad
 c_n=\frac{(2n+1)(2n+2)!}{(4n+2)!},
\]
where $f_\ell(n)=F_{i\ell}$ when $n=n_i$.
Apply these monic row combinations to
$H+U^{\mathrm{pert}}$, where
\[
 U^{\mathrm{pert}}_{tk}
 =\mathbf1_{k<5}\mu_{t+k-4}
 \frac{(2t+1)!(2b-2k)!}{(10-2k)!(2b+2t-9)!}.
\]
The last five transformed rows vanish in columns $k\ge5$.
The complementary block is upper triangular with diagonal
$h_0,\ldots,h_{b-6}$.  Reverse the retained five columns and
divide row $i$ by $c_{n_i}$.  The unperturbed entries become $F$.
For the perturbation use the elementary Bernoulli identity
\[
 \sum_{j=0}^{N}\binom Nj B_j(N-j)^{\underline h}
       =N^{\underline h}B_{N-h}(1).
\]
Since $b\ge10$, all the retained indices $N-h$ are at least two;
only odd $h$ survive.  Restoring the omitted $j=1$ term supplies
one further copy of $F$, explaining the necessary factor $2F$
in \eqref{ten:calA}.  This proves
\begin{equation}\label{ten:multiplier}
 \det M_{5,b}=(-1)^{5b-15}
 \frac{\prod_{t=0}^{b-1}(2b+2t-9)!}
      {\prod_{k=0}^{b-1}(2b-2k)!}
 \left(\prod_{n=0}^{b-6}h_n\right)
 \left(\prod_{i=0}^{4}c_{n_i}\right)\det\mathcal A.
\end{equation}
The same argument with the perturbation absent gives the identical
multiplier for $\det H$ and $\det F$, proving \eqref{ten:rho}.

\subsection{One additional bit removes the remaining finite obstruction}

All congruences in this subsection are in $\mathbb Z_{(2)}$,
the rationals with odd denominator.  We use the classical fact
$v_2(B_{2m})=-1$ for $m\ge1$.

\begin{lemma}[A doubled Bernoulli congruence]\label{ten:B17}
For every integer $m\ge13$,
\[
 2B_{2m}\equiv\mathcal B_{17}(m)\pmod{131072},
\]
where
\begin{equation}\label{ten:Bpoly}
 \mathcal B_{17}(x)=1+20068x-26384x^2+2816x^3+45824x^4
             +4096x^5-24576x^6+65536x^8.
\end{equation}
\end{lemma}
\begin{proof}
Faulhaber's identity, with $L=512$, gives
\[
 \sum_{a=0}^{L-1}a^{2m}
 =LB_{2m}+\sum_{h=2}^{2m+1}
      \binom{2m}{h-1}\frac{L^h}{h}B_{2m+1-h}.
\]
The $h=2$ term vanishes.  Every term with $h\ge3$ has valuation
at least $9h-v_2(h)-1\ge26$.  Divide by $256$ and reduce modulo
$2^{17}$.  The even $a$ terms vanish because $2m\ge26$.
For odd $a$, expand
$a^{2m}=(1+(a^2-1))^m$; since $8\mid a^2-1$, terms with
binomial degree at least nine have valuation at least $27-8>17$
after division.  Therefore
\begin{equation}\label{ten:finite-average}
 2B_{2m}\equiv\frac1{256}
 \sum_{a=1,3,\ldots,511}\sum_{j=0}^{8}
          \binom mj(a^2-1)^j\pmod{2^{17}}.
\end{equation}
Expand the finite right side as a polynomial in $m$.  Its rational
coefficients all have odd denominators, and their reductions modulo
$2^{17}$, in increasing powers, are
\[
 (1,20068,-26384,2816,45824,4096,-24576,0,65536).
\]
This is an exact finite coefficient calculation, reproduced by the
verifier from the $256$ odd support points and the nine displayed
binomial terms.  It gives \eqref{ten:Bpoly} and proves the lemma.
\end{proof}

Define
\begin{equation}\label{ten:Rpoly}
 \begin{split}
 \mathcal R(b)={}&1536b^8+128b^7-928b^6-816b^5+464b^4\\
                 &-1696b^3-1930b^2-1831b+694.
 \end{split}
\end{equation}
The new coefficientwise polynomial identity is
\begin{equation}\label{ten:congruence}
 P_{10}\bigl(b;\mathcal B_{17}(b-1),\ldots,
                    \mathcal B_{17}(b-9)\bigr)
       \equiv32\mathcal R(b)\pmod{131072}.
\end{equation}
This is verified in $\mathbb Z[b]$, not only at selected integer
arguments.  A compact exact procedure is to group the numerator's
$10403$ terms by their $439$ $U$-monomials, substitute the nine
explicit polynomials, multiply modulo $131072$, and compare every
coefficient with \eqref{ten:Rpoly}.  The provided script carries
out precisely those finite ring operations, after independently
rebuilding the numerator from \eqref{ten:calA}.  Thus
\eqref{ten:congruence} is a finite algebraic certificate whose
entire input is displayed or defined by finite products above.

The linear coefficient of $\mathcal R$ is odd, and every coefficient
of degree at least two is even.  Hence for distinct integers $x,y$,
\begin{equation}\label{ten:isometry}
 v_2(\mathcal R(x)-\mathcal R(y))=v_2(x-y).
\end{equation}
Indeed, the quotient by $x-y$ is the odd integer $-1831$ plus an
even integer.  Direct substitution gives
\[
 \mathcal R(3010)\equiv0\pmod{4096},\qquad
 \mathcal R(962)\equiv2048\pmod{4096}.
\]
Equation \eqref{ten:isometry} proves that $3010$ is the unique
zero class modulo $4096$.

\begin{proposition}[The new uniform congruence range]\label{ten:range}
For $b\ge22$ and $b\not\equiv3010\pmod{4096}$,
$M_{5,b}$ is nonsingular, and
\begin{equation}\label{ten:valuation}
 v_2(\rho_{5,b})=v_2(b-3010)-5.
\end{equation}
\end{proposition}
\begin{proof}
Every $b-j$ in \eqref{ten:rho} is at least $13$.  Thus
\eqref{ten:B17} and \eqref{ten:congruence} give
\[
 \frac{P_{10}(b;2B_{2b-2},\ldots,2B_{2b-18})}{32}
          \equiv\mathcal R(b)\pmod{4096}.
\]
Outside the unique zero class, \eqref{ten:isometry} gives
$v_2(\mathcal R(b))=v_2(b-3010)<12$.
The numerator therefore has valuation $5+v_2(b-3010)$.
The denominator in \eqref{ten:rho} has valuation $10$:
$138240=2^{10}\cdot135$, and all three remaining factors are odd
at integer $b$.  This proves the nonvanishing and
\eqref{ten:valuation}.
\end{proof}

\subsection{A self-contained sufficient real-variable threshold}

We specialize the general highest-degree argument of
\cite{QuarticIncoming}; this subsection is a credited reproduction
with the explicit order-ten constants independently recomputed.
Let
\[
 J(b)=\prod_{i=0}^{4}(2n_i+1)(2n_i+2),\qquad
 E(b)=\prod_{i,\ell=0}^{4}\prod_{r=1}^{\max\{0,2i-4\ell-1\}}
                  (R_\ell+r),\qquad D(b)=10!J(b)E(b).
\]
Here $\deg E=20$ and $\deg D=30$.  Clearing the fixed factorial
and row denominators in \eqref{ten:calA} gives
\[
 \widetilde{\mathcal A}=D\mathcal A
     \in\mathbb Z[b,U_1,\ldots,U_9]^{5\times5},\qquad
 \mathscr P=\det\widetilde{\mathcal A}.
\]
Every entry has $b$-degree at most $46$, so $\deg_b\mathscr P\le230$.
If $\|f\|_1$ denotes the sum of absolute values of all coefficients,
exact expansion of the $25$ entries gives
\[
 E^*:=\max_{i,\ell}\|\widetilde{\mathcal A}_{i\ell}\|_1
 =602847821074996042728613866361551228764160.
\]
Consequently $\|\mathscr P\|_1\le K:=120(E^*)^5$.
The number $E^*$ is independently recomputed from the entry formulas
by the accompanying threshold script; it is not inferred from
numerical samples.

Put $F_0=(2b-18)!$ and
\[
 V_j=U_j\frac{F_0}{(2b-2j)!},\qquad
 \kappa=\frac{120(10!)^5}{2!4!6!8!10!}
         =14932726824960000000.
\]
The part of $\mathscr P$ homogeneous of total degree five in $U$
has the exact form
\begin{equation}\label{ten:leading}
 \mathscr P_{[5]}=-\mathscr G(b)\det(V_{i+j+1})_{0\le i,j\le4},
\end{equation}
where
\[
 \mathscr G(b)=\kappa J(b)^4E(b)^5
 \prod_{i=0}^{4}(4n_i+2)^{\underline{2i}}
 \prod_{\ell=0}^{4}\prod_{r=1}^{10+2\ell}(2b-18+r).
\]
To see the structural factorization, write
$T=\mathcal A-2F$ and use $U_j=(2b-2j)!V_j/F_0$ in
\eqref{ten:calA}.  The identity
$N_{i\ell}-(2k+1)=2b-2(5-i+\ell+k)$ cancels the factorials.
The remaining matrix factors as
\[
 T=\operatorname{diag}\!\left(\frac1{(2n_i+1)(2n_i+2)}\right)
       L\,(V_{5-i+\ell})_{0\le i,\ell\le4}
       \operatorname{diag}\!\left(\frac{R_\ell!}{F_0(2\ell+2)!}\right),
\]
where $L_{i,i-k}=d_{i,2k+1}$ is lower triangular.  Its diagonal is
$d_{i1}=-(i+1)(4n_i+2)^{\underline{2i}}$, which follows directly
from the difference of \eqref{ten:Newton} at $y=1$ and $y=0$.
Reverse the five middle rows and multiply all row factors and
$D^5$.  This gives exactly \eqref{ten:leading}.

At $U_j=2B_{2b-2j}$, Euler's Bernoulli formula yields
\begin{equation}\label{ten:actualleading}
 |\mathscr P_{[5]}|
 =\frac{\mathscr G(b)(4F_0)^5}{(2\pi)^{10b-50}}
       \det\bigl(\zeta(2b-18+2i+2j)\bigr)_{0\le i,j\le4}.
\end{equation}
The determinant is strictly positive.  Indeed, Cauchy--Binet for
the discrete moments on $m^{-2}$ gives
\[
 \det\bigl(\zeta(2b-18+2i+2j)\bigr)
 =\sum_{m_1<\cdots<m_5}(m_1\cdots m_5)^{-2b+18}
                  \prod_{r<s}(m_r^{-2}-m_s^{-2})^2.
\]
Keeping $(m_1,\ldots,m_5)=(1,2,3,4,5)$ gives the lower bound
\[
 V\,120^{-2b+18},\qquad
 V=\prod_{1\le r<s\le5}(r^{-2}-s^{-2})^2
          =\frac{2401}{36864000000000000}.
\]
For $b\ge20$ every linear factor of $\mathscr G$ is at least $b$,
and its degree is $230$, so $\mathscr G(b)\ge\kappa b^{230}$.
For $b\ge30$,
\[
 |2B_{2b-2j}|\le
 T_b:=\frac{8(2b)^{16}(2b-18)!}{(2\pi)^{2b-18}},
 \qquad 1\le j\le9,
\]
using $\zeta(2b-2j)<2$.  Also $T_b\ge1$: since $2b-18\ge42$,
the elementary comparison $26!>8^{26}$ and induction suffice.
All lower-degree terms of $\mathscr P$ have total $U$-degree at
most four.  Therefore
\[
 |\mathscr P_{[\le4]}|\le Kb^{230}T_b^4.
\]
Divide by \eqref{ten:actualleading}, cancel the powers of $2\pi$
before using $2\pi<8$, and obtain
\begin{equation}\label{ten:bound}
 \frac{|\mathscr P_{[\le4]}|}{|\mathscr P_{[5]}|}
 \le R(b):=C\frac{b^{64}960^{2b}}{(2b-18)!},\qquad
 C=\frac{K8^{26}2^{64}}{\kappa4^5V120^{18}}.
\end{equation}
Every constant is explicit and rational.  Exact integer comparison
after clearing denominators gives
\[
 R(1920)<2^{-546}<1.
\]
The verifier checks this stronger margin.  Moreover, for $b\ge1920$,
\[
 \frac{R(b+1)}{R(b)}=
 \left(1+\frac1b\right)^{64}
      \frac{960^2}{(2b-16)(2b-17)}<\frac34.
\]
For a simple uniform justification, $(1+1/b)^{64}<3$,
both denominator factors exceed $b$, and $b\ge1920=2\cdot960$.
Thus the highest-degree term cannot be cancelled for any $b\ge1920$.
Since $D(b)>0$ and the multiplier \eqref{ten:multiplier} is nonzero,
we have proved
\begin{equation}\label{ten:eventual}
             \det M_{5,b}\ne0\qquad(b\ge1920).
\end{equation}

\subsection{Completion, and an independent audit of the old exception}

For $22\le b<1920$, one cannot have $b\equiv3010\pmod{4096}$.
Proposition~\ref{ten:range} therefore proves nonvanishing throughout
this interval.  Equation \eqref{ten:eventual} supplies all
$b\ge1920$.  The remaining $6\le b\le21$ are checked directly
from the rational source matrix \eqref{ten:M}; the machine-readable
record contains every nonzero determinant.  Most of these base cases
are already contained in \cite{AlternatingIncoming}; they are recomputed here.
For example,
\[
 \det M_{5,6}=-\frac{7097067}{2867200},\qquad
 \det M_{5,8}=\frac{537653622905129}{68669669376000}.
\]
This proves Theorem~\ref{ten:main}.

The old exceptional index $b=962$ also has a short independent exact
audit.  Its numerator satisfies
\begin{equation}\label{ten:962}
 P_{10}(962;2B_{1922},2B_{1920},\ldots,2B_{1906})
       \equiv2386755584\pmod{2^{32}}.
\end{equation}
The residue has valuation $16$, giving $v_2(\rho_{5,962})=6$.
Two independent calculations of the nine Bernoulli residues agree:
the recurrence
\[
 (2m+1)(2B_{2m})=2m-1-
        \sum_{j=1}^{m-1}\binom{2m+1}{2j}(2B_{2j}),
\]
with odd division modulo $2^{32}$, and the Faulhaber power sum over
the $32768$ odd integers below $65536$, performed modulo $2^{47}$
before dividing by $2^{15}$.  Finally, evaluating the rational
five-dimensional determinant \eqref{ten:calA} directly, with
these integer residues as inputs, gives the same residue
\eqref{ten:962} as the sparse numerator.  This audit is not needed
once the uniform congruence is proved, but it verifies the exact
location where the two preceding reports failed to overlap.

\subsection{An ordinary convergent family of polygamma identities}

For every $b\ge6$ and $0\le k<b$, let
\begin{equation}\label{ten:Q}
 Q_{b,k}(x)=\sum_{t=0}^{b-1}(M_{5,b}^{-1})_{kt}
       \bigl((x+\tfrac12)^{2b+2t+1}
                    -(x-\tfrac12)^{2b+2t+1}\bigr).
\end{equation}
This form is shorter than an expansion of the even polynomial.
The identity
\[
 (x+\tfrac12)^{2j+1}-(x-\tfrac12)^{2j+1}
 =\sum_{r=0}^{j}\binom{2j+1}{2r}2^{-2(j-r)}x^{2r}
\]
shows that \eqref{ten:Q} is precisely the inverse coordinate
combination dictated by \eqref{ten:C}.

Define
\[
 \eta_{s,0}=\frac1{s-1},\qquad
 \eta_{s,j}=\frac{(s)_{2j-1}B_{2j}(1/2)}{(2j)!}\quad(j\ge1),
 \qquad c_{b,\ell}=\sum_{j=0}^{\ell}\eta_{10,j}\eta_{2b,\ell-j}.
\]
Let $S_{b,k}(x)$ be the polynomial consisting of the nonnegative
powers in the formal expansion
\[
 Q_{b,k}(x)x^{-2b-8}\sum_{\ell\ge0}c_{b,\ell}x^{-2\ell}.
\]
It is an even polynomial of degree at most $2b-10$, with rational
coefficients determined by finitely many terms.

\begin{corollary}[All order-ten reconstruction identities]
\label{ten:identity}
For every integer $b\ge6$ and $0\le k<b$,
\begin{equation}\label{ten:polygamma}
 \boxed{\displaystyle
 \sum_{n=0}^{\infty}\left\{
 Q_{b,k}(x_n)
 \frac{\psi^{(9)}(n+1)\psi^{(2b-1)}(n+1)}{9!\,(2b-1)!}
             -S_{b,k}(x_n)\right\}=\zeta(2k).}
\end{equation}
The series is ordinary and absolutely convergent.
\end{corollary}
\begin{proof}
For even $s$, $A_s(n)=\psi^{(s-1)}(n+1)/(s-1)!$.
By \eqref{ten:C}, \eqref{ten:Q}, and the exact matrix inverse,
the continued sum of $Q_{b,k}(x_n)A_{10}(n)A_{2b}(n)$ is
$\zeta(2k)$.  The centered Euler--Maclaurin expansion of the product
is exactly the formal product used to define $S_{b,k}$.
Removing all nonnegative powers leaves $O(x_n^{-2})$.
Every removed power is even and has centered continued sum
$\zeta(-2r,1/2)=0$, $r\ge0$.  Therefore the continued sum agrees
with the displayed absolutely convergent series.
\end{proof}

Here is a completely explicit member, included to make the
reconstruction mechanism concrete.  Let
\begin{align*}
 P(x)={}&58535757x^{23}+145331415x^{21}-98765059x^{19}\\
        &+92514527x^{17}-73310041x^{15}+28468113x^{13}.
\end{align*}
Then
\begin{equation}\label{ten:explicitidentity}
 \begin{split}
 \sum_{n=0}^{\infty}\Bigg\{&
    [P(n+1)-P(n)]
        [\zeta(10)-H_n^{(10)}][\zeta(12)-H_n^{(12)}]\\
   &-\frac{149591379(n+\tfrac12)^2+1835020425}{11}
                          \Bigg\}
       =-\frac{574862427}{8}.
 \end{split}
\end{equation}
It follows from the $k=0$ row of $M_{5,6}^{-1}$:
\[
 \frac4{574862427}
 (28468113,-73310041,92514527,-98765059,145331415,58535757).
\]
The numerical pair $(10,12)$ was already within the source's finite
range; the explicit identity is an illustration of the now-uniform
construction, not the new nonvanishing case.  All coefficients,
counterterms, and the exact inverse-row check are included in the
identity artifact.

\paragraph{Further questions specific to this completion.}
The unrestricted rank problem remains open in the inspected corpus
for smaller order at least twelve.  A natural next question is whether
one can prove a general overlap principle between the finite
$2$-adic exceptional classes and the constructive real-variable
threshold, without computing an increasingly large numerator.
Another is whether a determinant representation makes the lower
Bernoulli-degree terms sign-compatible, giving nonvanishing directly
and avoiding a domination threshold.  The gap-six family, where the
generic nearby-Bernoulli contraction falls outside its stable range,
still calls for a separate argument.  Finally, one may seek compact
closed forms for the inverse rows in \eqref{ten:Q}; the existence
and rational computability proved here do not by themselves make
those coefficients small.
